\documentclass[12pt,oneside]{book}

% ==========================================================
% METADATOS
% ==========================================================
\newcommand{\universidad}{Universidad de Buenos Aires}
\newcommand{\facultad}{Facultad de Derecho}
\newcommand{\materia}{Derecho Procesal Civil y Comercial}
\newcommand{\titulo}{Notificaciones, Rebeld\'ia y Pr\'actica Profesional en el Proceso de Familia}
\newcommand{\autor}{Isaac Edgar Camacho Ocampo}
\newcommand{\anio}{2026}

% ==========================================================
% PAQUETES
% ==========================================================
\usepackage[utf8]{inputenc}
\usepackage[T1]{fontenc}
\usepackage[spanish,es-nodecimaldot]{babel}

\usepackage[
    top=2.5cm,
    bottom=2.5cm,
    left=3cm,
    right=2.5cm,
    headheight=15pt
]{geometry}

\usepackage{amsmath}
\usepackage{amssymb}
\usepackage{mathtools}

\usepackage{graphicx}
\usepackage{float}
\usepackage{caption}
\usepackage{subcaption}

\usepackage{booktabs}
\usepackage{array}
\usepackage{longtable}
\usepackage{tabularx}

\usepackage{enumitem}

\usepackage{xcolor}
\usepackage[most]{tcolorbox}

\usepackage{fancyhdr}
\usepackage{ragged2e}

\usepackage[
    colorlinks=true,
    linkcolor=blue!50!black,
    citecolor=green!40!black,
    urlcolor=blue!60!black,
    pdftitle={\titulo},
    pdfauthor={\autor},
    pdfsubject={Apuntes universitarios},
    bookmarks=true
]{hyperref}

% ==========================================================
% CONFIGURACI\'ON
% ==========================================================
\graphicspath{
    {./img/}
    {./graficos/}
    {./figuras/}
}

\setcounter{tocdepth}{2}

\setlength{\parindent}{0pt}
\setlength{\parskip}{0.5em}

\setlist[itemize]{
    topsep=0.3em,
    itemsep=0.2em
}

\setlist[enumerate]{
    topsep=0.3em,
    itemsep=0.2em
}

\clubpenalty=10000
\widowpenalty=10000

% ==========================================================
% COMANDOS
% ==========================================================
\newcommand{\abs}[1]{\left|#1\right|}
\newcommand{\norm}[1]{\left\lVert#1\right\rVert}
\newcommand{\vect}[1]{\mathbf{#1}}
\newcommand{\deriv}[2]{\frac{d#1}{d#2}}
\newcommand{\pderiv}[2]{\frac{\partial#1}{\partial#2}}

% ==========================================================
% ENTORNOS / CAJAS
% ==========================================================
\newtcolorbox{definition}[1]{
    enhanced,
    breakable,
    title={Definici\'on: #1},
    fonttitle=\bfseries,
    colback=blue!3,
    colframe=blue!50!black,
    boxrule=0.8pt,
    arc=2mm
}

\newtcolorbox{important}{
    enhanced,
    breakable,
    title={Importante},
    fonttitle=\bfseries,
    colback=yellow!8,
    colframe=orange!70!black,
    boxrule=0.8pt,
    arc=2mm
}

\newtcolorbox{example}[1]{
    enhanced,
    breakable,
    title={Ejemplo: #1},
    fonttitle=\bfseries,
    colback=green!4,
    colframe=green!40!black,
    boxrule=0.8pt,
    arc=2mm
}

\newtcolorbox{formula}[1]{
    enhanced,
    breakable,
    title={F\'ormula / Regla: #1},
    fonttitle=\bfseries,
    colback=purple!4,
    colframe=purple!50!black,
    boxrule=0.8pt,
    arc=2mm
}

\newtcolorbox{exercise}[1]{
    enhanced,
    breakable,
    title={Ejercicio: #1},
    fonttitle=\bfseries,
    colback=gray!5,
    colframe=gray!60!black,
    boxrule=0.8pt,
    arc=2mm
}

\newtcolorbox{note}{
    enhanced,
    breakable,
    title={Nota},
    fonttitle=\bfseries,
    colback=white,
    colframe=gray!50,
    boxrule=0.6pt,
    arc=2mm
}

% ==========================================================
% CONFIGURACI\'ON FINAL
% ==========================================================
\pagestyle{fancy}
\fancyhf{}
\fancyhead[L]{\nouppercase{\leftmark}}
\fancyhead[R]{\materia}
\fancyfoot[C]{\thepage}
\renewcommand{\headrulewidth}{0.4pt}
\renewcommand{\footrulewidth}{0pt}

\fancypagestyle{plain}{
    \fancyhf{}
    \fancyfoot[C]{\thepage}
    \renewcommand{\headrulewidth}{0pt}
}

% Anchos del cuadro Cornell (columna izquierda estrecha, derecha amplia)
\newcommand{\cornellL}{0.27\textwidth}
\newcommand{\cornellR}{0.63\textwidth}

\begin{document}

% ==========================================================
% PORTADA
% ==========================================================
\begin{titlepage}

\thispagestyle{empty}

\begin{center}

\vspace*{1cm}

{\large \textsc{\universidad}}

\vspace{0.2cm}

{\large \textsc{\facultad}}

\vspace{3cm}

{\Huge \textbf{\materia}}

\vspace{1cm}

{\LARGE \titulo}

\vspace{0.8cm}

\textit{Comprender los conceptos es el primer paso para convertir el estudio en conocimiento.}

\vspace{1.2cm}

\rule{0.70\textwidth}{0.5pt}

\vspace{0.5cm}

{\large Apuntes de estudio}

\vspace{0.5cm}

\rule{0.70\textwidth}{0.5pt}

\vfill

{\large \autor}

\vspace{0.3cm}

{\large \anio}

\end{center}

\vfill

\begin{flushleft}
\textit{Este es un modesto aporte a los estudiantes de ``Derecho Procesal Civil y Comercial'' no pretende sustituir el material oficial de la materia.}
\end{flushleft}

\end{titlepage}

% ==========================================================
% \'INDICE
% ==========================================================
\tableofcontents

% ==========================================================
% CONTENIDO
% ==========================================================
\chapter{Presentaci\'on de la clase}
\chaptermark{Presentaci\'on}

\section{Temas que aborda la clase}

La clase combina dos momentos de una misma pr\'actica procesal. En primer lugar, un repaso te\'orico sobre el tr\'amite de la demanda, el traslado, el \textbf{principio de preclusi\'on}, el r\'egimen de \textbf{notificaciones} (art\'iculo 135 del C\'odigo Procesal Civil y Comercial de la Naci\'on, CPCCN) y la declaraci\'on de \textbf{rebeld\'ia} (art\'iculo 59), con su importante excepci\'on en los procesos de alimentos. En segundo lugar, un taller pr\'actico en el que se redacta y corrige un escrito judicial que solicita una medida urgente de vinculaci\'on provisoria dentro de un proceso de familia, trabajando sobre la constituci\'on de domicilio y la redacci\'on del petitorio.

Los contenidos se organizan en tres bloques:
\begin{itemize}
    \item \textbf{Bloque I.} Marco te\'orico procesal: la demanda, el traslado y las notificaciones.
    \item \textbf{Bloque II.} La rebeld\'ia y su excepci\'on en el derecho de familia.
    \item \textbf{Bloque III.} Taller pr\'actico de redacci\'on: medida de vinculaci\'on provisoria.
\end{itemize}

Los tres bloques est\'an unidos por una misma idea: el camino que recorre un expediente judicial desde que se presenta la demanda hasta que se convierte en un escrito real, con domicilios, plazos y pedidos concretos.

\begin{note}
Los apuntes originales presentan este recorrido mediante una analog\'ia con un viaje en auto: primero hay que conocer las reglas de tr\'ansito (la demanda, el traslado, la preclusi\'on); despu\'es, entender c\'omo se avisa a los dem\'as conductores lo que uno va a hacer (las notificaciones y la rebeld\'ia); y, finalmente, manejar de verdad, corrigiendo el rumbo sobre la marcha (el taller de redacci\'on del escrito de vinculaci\'on provisoria).
\end{note}

\section{Utilidad para la vida profesional}

Seg\'un los apuntes, todo abogado litigante, sin importar el fuero en el que trabaje, necesita dominar estos contenidos desde el primer d\'ia de ejercicio profesional:
\begin{itemize}
    \item Un error en el c\'omputo del traslado o en la forma de notificar puede provocar la \textbf{nulidad} de todo lo actuado y hacer perder un derecho a un cliente real.
    \item Comprender el principio de preclusi\'on evita presentar escritos fuera de t\'ermino o insistir sobre etapas ya cerradas del proceso.
    \item Saber cu\'ando corresponde pedir la rebeld\'ia, y cu\'ando no (como en los procesos de alimentos), evita cometer un error grave frente a un juez y frente al cliente.
    \item El taller de redacci\'on muestra c\'omo se corrige un escrito bajo presi\'on de tiempo: revisar domicilios, ajustar el petitorio y coordinar la presentaci\'on electr\'onica son tareas cotidianas en cualquier estudio jur\'idico.
\end{itemize}

\begin{important}
Esta clase es la caja de herramientas m\'inima para no perder un juicio por un error de forma y para actuar con responsabilidad ante situaciones urgentes que afectan a personas reales, como ni\~nas, ni\~nos y adolescentes.
\end{important}

% ----------------------------------------------------------
\chapter{Marco te\'orico procesal: la demanda, el traslado y las notificaciones}
\chaptermark{Marco te\'orico procesal}

\section{Repaso: la demanda y su tr\'amite inicial}

La clase retoma el tema desarrollado anteriormente: c\'omo se inicia una demanda y cu\'al es su circuito. Se remarca un aspecto central de ese circuito: el control de los plazos y la obligaci\'on de \textbf{verificar el estado del expediente dentro de las veinticuatro horas}; si no se lo verifica en ese lapso, debe darse aviso. Se lo se\~nala como casi lo m\'as importante de lo visto.

\section{La demanda, el traslado y el primer prove\'ido}

Una vez que la demanda est\'a presentada en el juzgado, el juez, al recibirla, dicta el \textbf{primer prove\'ido}.

\subsection{Funci\'on del primer prove\'ido}

El primer prove\'ido \textbf{ordena el traslado}. En \'el, el juez determina qu\'e hay, qu\'e no hay, qu\'e falta y qu\'e no falta. Si no existe nada que el juez necesite aclarar o exigir, ordena el traslado. Una vez superado ese control inicial (las partes, el objeto, el hecho y el derecho est\'an correctamente expuestos), corresponde el traslado.

\subsection{El prove\'ido ``previo''}

Cada prove\'ido que dice \textbf{``previo''} significa que, antes de pasar a la otra parte, debe cumplirse algo; si no se cumple, el tr\'amite no avanza. Funciona, as\'i, como un control de admisibilidad antes de continuar el tr\'amite.

\section{El principio de preclusi\'on}

A partir de la idea del ``previo'', se desarrolla uno de los principios estructurales del proceso civil.

\begin{definition}{Principio de preclusi\'on}
Las etapas del proceso civil se suceden en un orden fijo y, una vez cerrada una etapa (por ejemplo, el plazo para contestar la demanda), la parte pierde definitivamente la posibilidad de realizar el acto que omiti\'o realizar en su momento. El proceso civil es un sistema ``estanco'': \textbf{va hacia adelante y no se vuelve}.
\end{definition}

Lo que no se hizo en la etapa correspondiente no puede realizarse despu\'es, salvo las excepciones que la propia ley contempla. Si del otro lado hay una parte atenta, har\'a saber que la etapa fue superada y que se perdi\'o un derecho.

\section{Las notificaciones y el art\'iculo 135 del CPCCN}

El art\'iculo 135 del CPCCN es el que establece las formas de notificaci\'on. Se indica su lectura, especialmente en lo relativo a las c\'edulas (que son varias), y se se\~nala que no dice mucho m\'as porque el C\'odigo es muy antiguo. Posteriormente se analizar\'a la distinci\'on entre notificaci\'on electr\'onica y notificaci\'on en papel.

% TODO: contenido incompleto o ambiguo en las notas originales (el texto del art. 135 no fue transcripto; la caja siguiente es una s\'intesis elaborada en los apuntes originales)
\begin{important}
\textbf{Art\'iculo 135 del CPCCN (notificaci\'on por c\'edula).} El art\'iculo enumera de modo taxativo los supuestos en los que las resoluciones judiciales deben notificarse por c\'edula (entre ellos, el traslado de la demanda y de la reconvenci\'on, la citaci\'on de personas ajenas al proceso, la sentencia definitiva y, en general, toda resoluci\'on que la ley mande notificar de esa forma expresamente).

Los apuntes aclaran que se trata de una s\'intesis elaborada a partir de lo se\~nalado en clase y recomiendan cotejar el texto vigente en una edici\'on comentada del C\'odigo.
\end{important}

\subsection{La necesidad de un C\'odigo comentado}

Se insiste en que es fundamental contar con un \textbf{c\'odigo comentado}, tanto del C\'odigo Civil y Comercial como del C\'odigo Procesal, en el formato que cada estudiante prefiera, porque ser\'an de uso constante. Se remarca que la tecnolog\'ia no reemplaza el criterio jur\'idico.

\begin{note}
La biblioteca digital de la instituci\'on da acceso a dos bases de datos jur\'idicas, y en alguna de ellas hay comentario de cada art\'iculo; tambi\'en se menciona la biblioteca de la facultad. % TODO: contenido incompleto o ambiguo en las notas originales (las bases de datos se mencionan de modo coloquial, sin nombre preciso)
\end{note}

\section{La c\'edula electr\'onica en el expediente digital}

Se indica leer el art\'iculo 135 y observar qu\'e corresponde a la notificaci\'on electr\'onica y qu\'e a la notificaci\'on en papel. Tambi\'en se se\~nala d\'onde verificar la c\'edula electr\'onica en el sistema de gesti\'on judicial (\textbf{PJN}): en cada uno de los expedientes, en especial los reci\'en iniciados, hay una c\'edula electr\'onica.

\begin{important}
En los expedientes tramitados electr\'onicamente, la c\'edula de notificaci\'on se genera y se agrega en formato digital dentro del propio expediente virtual (sistema PJN), pudiendo verificarse su diligenciamiento en cualquiera de las causas recientemente iniciadas. Este mecanismo se distingue del r\'egimen de la c\'edula en soporte papel, cuyo tratamiento se remite a una clase posterior.
\end{important}

\subsection*{Cuadro Cornell: Bloque I}
\addcontentsline{toc}{section}{Cuadro Cornell: Bloque I}

\begin{longtable}{@{}>{\raggedright\arraybackslash}p{\cornellL}@{\hspace{1em}}>{\raggedright\arraybackslash}p{\cornellR}@{}}
\toprule
\textbf{Preguntas / palabras clave} & \textbf{Notas / respuestas} \\
\midrule
\endhead
\bottomrule
\endfoot

\textbf{¿Qu\'e hace el juez al recibir la demanda y qu\'e ordena el primer prove\'ido?} &
El juez dicta el primer prove\'ido, que ordena el traslado, luego de revisar que est\'en bien el objeto, el hecho y el derecho. En ese prove\'ido determina qu\'e hay, qu\'e no hay, qu\'e falta y qu\'e no falta. \\[0.8em]

\textbf{¿Qu\'e significa que un prove\'ido diga ``previo''?} &
Que, antes de pasar a la otra parte, debe cumplirse algo; si no, el tr\'amite no avanza. \\[0.8em]

\textbf{¿Qu\'e es el principio de preclusi\'on?} &
Es el principio por el cual las etapas del proceso se suceden en un orden fijo e irreversible: una vez cerrada una etapa procesal, no se puede volver a ella, y lo que no se hizo en el momento oportuno se pierde como derecho, salvo excepci\'on legal expresa. \\[0.8em]

\textbf{¿Qu\'e regula el art\'iculo 135 del CPCCN?} &
Las formas de notificaci\'on de las resoluciones judiciales, en especial la notificaci\'on por c\'edula, que corresponde para el traslado de la demanda y otras resoluciones de particular importancia para el ejercicio del derecho de defensa. \\[0.8em]

\textbf{¿D\'onde se verifica una notificaci\'on electr\'onica?} &
En los expedientes que tramitan electr\'onicamente (sistema PJN), la c\'edula se genera y consulta dentro del propio expediente digital, disponible desde el inicio de la causa. \\

\midrule
\multicolumn{2}{@{}p{\dimexpr\cornellL+\cornellR+1em\relax}@{}}{%
\textbf{Resumen:} El primer prove\'ido del juez ordena el traslado una vez controlados los requisitos; si falta algo, el prove\'ido ``previo'' exige cumplirlo antes de continuar. La preclusi\'on impide retroceder a etapas cerradas. El art\'iculo 135 del CPCCN regula las formas de notificaci\'on (en especial la c\'edula), y en los expedientes electr\'onicos la c\'edula se verifica en el propio expediente digital del PJN.} \\
\end{longtable}

% ----------------------------------------------------------
\chapter{La rebeld\'ia y su excepci\'on en el derecho de familia}
\chaptermark{La rebeld\'ia}

\section{La declaraci\'on de rebeld\'ia (art\'iculo 59 del CPCCN)}

Se indica como lectura el art\'iculo 59 del CPCCN, descripto como muy simple y concreto. El supuesto de hecho que da lugar a la rebeld\'ia es el siguiente: quien \textbf{no contesta cuando debe contestar} ser\'a declarado rebelde.

La rebeld\'ia se declara mediante una \textbf{resoluci\'on del juez}, que debe ser notificada.

% TODO: contenido incompleto o ambiguo en las notas originales (la respuesta del alumnado sobre qu\'e hacer con la resoluci\'on no qued\'o registrada con claridad; la caja siguiente es una s\'intesis elaborada en los apuntes originales)
\begin{important}
\textbf{Art\'iculo 59 del CPCCN (declaraci\'on de rebeld\'ia).} Si la parte debidamente citada no compareciere durante el plazo de la citaci\'on, o abandonare el proceso despu\'es de haber comparecido, se la declarar\'a en rebeld\'ia a pedido de la otra parte. Esa resoluci\'on debe notificarse por c\'edula o, en su caso, por los medios que correspondan, precisamente para que quien resulta declarado rebelde tenga la posibilidad de comparecer en cualquier estado del proceso.

Los apuntes aclaran que se trata de una s\'intesis elaborada a partir de la explicaci\'on dada en clase y de la doctrina procesal habitual sobre la norma, y recomiendan cotejar el texto vigente en el C\'odigo comentado.
\end{important}

\section{La rebeld\'ia en el proceso de alimentos}

\begin{important}
\textbf{En el proceso de alimentos no hay rebeld\'ia.} Pedir la rebeld\'ia cuando el alimentante no contesta la demanda constituye un error grave.
\end{important}

\subsection{Fundamento: la funci\'on tuitiva del proceso de alimentos}

El proceso de familia de alimentos tiene una funci\'on pr\'actica: \textbf{busca que el alimentante pague siempre}. Declarar la rebeld\'ia implicar\'ia sacarlo del proceso por un tiempo, y el proceso nunca pretende excluirlo sino, por el contrario, mantenerlo dentro para que pague y se haga cargo de los alimentados.

Est\'a ``bien pensado'': m\'as all\'a de que el alimentante no cumpla, en cualquier momento del proceso puede ingresar, de la manera que fuera, al proceso de alimentos.

\begin{important}
A diferencia de la regla general del art\'iculo 59 del CPCCN, en el proceso de alimentos no corresponde declarar la rebeld\'ia del alimentante que no contesta la demanda. El fundamento reside en la funci\'on tuitiva del proceso de alimentos: su finalidad es asegurar que el alimentante permanezca vinculado a la causa y cumpla con su obligaci\'on, y no expulsarlo del proceso. Por ello, la ley permite que el alimentante se presente e intervenga en cualquier estado del tr\'amite, sin las restricciones propias de la rebeld\'ia.
\end{important}

\subsection*{Cuadro Cornell: Bloque II}
\addcontentsline{toc}{section}{Cuadro Cornell: Bloque II}

\begin{longtable}{@{}>{\raggedright\arraybackslash}p{\cornellL}@{\hspace{1em}}>{\raggedright\arraybackslash}p{\cornellR}@{}}
\toprule
\textbf{Preguntas / palabras clave} & \textbf{Notas / respuestas} \\
\midrule
\endhead
\bottomrule
\endfoot

\textbf{¿Qu\'e es la rebeld\'ia y c\'omo se declara (art\'iculo 59 del CPCCN)?} &
Es la situaci\'on procesal de la parte que, debidamente citada, no contesta la demanda o abandona el proceso. Se declara mediante una resoluci\'on del juez, que debe notificarse para que la parte pueda comparecer en cualquier momento posterior. \\[0.8em]

\textbf{¿Por qu\'e no se aplica la rebeld\'ia en el proceso de alimentos?} &
Porque el proceso de alimentos tiene una funci\'on tuitiva: busca que el alimentante permanezca dentro del proceso y cumpla su obligaci\'on, no que sea excluido de \'el. Por eso puede presentarse en cualquier estado del tr\'amite sin las restricciones propias de la rebeld\'ia. \\

\midrule
\multicolumn{2}{@{}p{\dimexpr\cornellL+\cornellR+1em\relax}@{}}{%
\textbf{Resumen:} Quien no contesta cuando debe hacerlo es declarado rebelde mediante una resoluci\'on judicial que debe notificarse (art\'iculo 59 del CPCCN). En alimentos no hay rebeld\'ia: pedirla es un error grave, porque el proceso busca mantener al alimentante dentro de la causa para que pague, y \'el puede ingresar al proceso en cualquier momento.} \\
\end{longtable}

% ----------------------------------------------------------
\chapter{Taller pr\'actico de redacci\'on: medida de vinculaci\'on provisoria}
\chaptermark{Taller pr\'actico}

El segundo momento de la clase es eminentemente pr\'actico: consiste en la correcci\'on, en tiempo real y en grupo, de un escrito judicial correspondiente a un proceso de familia, mediante el cual se solicita una \textbf{medida urgente de vinculaci\'on provisoria}. El docente y los estudiantes trabajan, a trav\'es de una plataforma de edici\'on compartida, sobre el encabezado, el domicilio y el petitorio del escrito.

\section{Constituci\'on y denuncia del domicilio en el escrito}

La actividad comienza con la revisi\'on de un escrito ya redactado, al que es necesario incorporarle un \textbf{domicilio real} que hab\'ia sido omitido. La urgencia del pedido se vincula con una situaci\'on de temor manifestada por una de las partes: si ella dijo que ten\'ia miedo, se solicita la vinculaci\'on provisoria.

Del trabajo realizado surgen las siguientes observaciones:
\begin{itemize}
    \item Ese dato debe incorporarse formalmente en el \textbf{petitorio} del escrito.
    \item Tambi\'en corresponde modificar el \textbf{encabezado}, que es donde se ubica el domicilio real.
    \item Se incorpora la direcci\'on real denunciada: Moreno 2472, San Miguel. Se corrige un error de localidad: hab\'ia una confusi\'on con otra localidad, por lo que se aclara la ciudad. La correcci\'on se limita a precisar la ciudad, sin alterar el resto de los datos.
    \item La secci\'on se titula \textbf{``se constituye domicilio real''}.
\end{itemize}

\begin{important}
El taller distingue, en la pr\'actica, entre el \textbf{domicilio real} de la parte (el lugar f\'isico donde efectivamente reside; en este caso, Moreno 2472, San Miguel, provincia de Buenos Aires) y el \textbf{domicilio constituido} a los efectos procesales dentro del expediente. La correcta identificaci\'on de la ciudad o partido resulta indispensable para que el domicilio denunciado sea v\'alido y produzca los efectos procesales que la parte pretende (en este caso, sustentar el pedido de vinculaci\'on provisoria por razones de urgencia).
\end{important}

\section{El petitorio y la solicitud de vinculaci\'on provisoria}

Con el domicilio ya corregido, se determina si basta con la incorporaci\'on realizada o si hace falta algo m\'as: es necesario \textbf{modificar el petitorio}, y no agregar otro.

Se redacta entonces el nuevo punto del petitorio referido a la vinculaci\'on provisoria. En la redacci\'on colectiva se discuten las f\'ormulas posibles: no corresponde usar ``tenga'', sino que se pide que \textbf{se ordene la vinculaci\'on}. Se aclara tambi\'en que lo relativo a la constituci\'on del domicilio real es un punto distinto del pedido de vinculaci\'on.

\begin{formula}{Redacci\'on acordada del punto del petitorio}
``Se ordene la vinculaci\'on provisoria solicitada en el punto\ldots'' % TODO: contenido incompleto o ambiguo en las notas originales (la redacci\'on definitiva no qued\'o transcripta completa)
\end{formula}

El punto se ubica en el lugar correspondiente dentro del escrito: seg\'un los apuntes, en el \textbf{punto cuatro}. % TODO: contenido incompleto o ambiguo en las notas originales (el docente dice ``en el punto cuatro, creo'')

\begin{important}
Cuando se incorpora un pedido nuevo a un escrito ya redactado (en este caso, la vinculaci\'on provisoria fundada en el temor manifestado por una de las partes), no basta con mencionarlo en los considerandos: es necesario trasladar ese pedido, con redacci\'on precisa e identificando el punto correspondiente, a la parte petitoria del escrito, que es la secci\'on que efectivamente delimita lo que se solicita al juez que resuelva.
\end{important}

\section{Presentaci\'on del escrito}

El taller cierra con la coordinaci\'on de la presentaci\'on del escrito: se distribuyen las tareas de firma y carga del documento en el sistema, y se indica que debe subirse el mismo d\'ia. Tambi\'en quedan gestiones pendientes vinculadas con la causa.

\subsection*{Cuadro Cornell: Bloque III}
\addcontentsline{toc}{section}{Cuadro Cornell: Bloque III}

\begin{longtable}{@{}>{\raggedright\arraybackslash}p{\cornellL}@{\hspace{1em}}>{\raggedright\arraybackslash}p{\cornellR}@{}}
\toprule
\textbf{Preguntas / palabras clave} & \textbf{Notas / respuestas} \\
\midrule
\endhead
\bottomrule
\endfoot

\textbf{¿Qu\'e es la vinculaci\'on provisoria solicitada en el taller pr\'actico?} &
Es una medida urgente pedida dentro de un proceso de familia ante una situaci\'on de temor manifestada por una de las partes, orientada a resolver de manera transitoria el contacto entre las personas involucradas mientras se sustancia el proceso principal. \\[0.8em]

\textbf{¿Qu\'e diferencia hay entre domicilio real y domicilio constituido?} &
El domicilio real es el lugar f\'isico donde efectivamente vive la persona (en el caso trabajado en clase, Moreno 2472, San Miguel, provincia de Buenos Aires); el domicilio constituido es el que se fija a los efectos de las notificaciones dentro del expediente. Ambos deben quedar correctamente identificados, incluyendo la localidad, para que produzcan efectos v\'alidos. \\[0.8em]

\textbf{¿C\'omo se incorpora un pedido nuevo al petitorio de un escrito?} &
No alcanza con mencionarlo en el cuerpo del escrito: debe redactarse con precisi\'on dentro de la secci\'on de petitorio, identificando el punto correspondiente, ya que es esa secci\'on la que delimita exactamente lo que se solicita al juez que resuelva. \\

\midrule
\multicolumn{2}{@{}p{\dimexpr\cornellL+\cornellR+1em\relax}@{}}{%
\textbf{Resumen:} En el taller se corrigi\'o un escrito de familia que solicita una vinculaci\'on provisoria urgente: se incorpor\'o el domicilio real con su ciudad, se distingui\'o del domicilio constituido y se traslad\'o el nuevo pedido al petitorio, con redacci\'on precisa y ubicado en su punto.} \\
\end{longtable}

% ----------------------------------------------------------
\chapter{Resumen Cornell integrador}
\chaptermark{Resumen integrador}

\begin{longtable}{@{}>{\raggedright\arraybackslash}p{\cornellL}@{\hspace{1em}}>{\raggedright\arraybackslash}p{\cornellR}@{}}
\toprule
\textbf{Preguntas / palabras clave} & \textbf{Notas / respuestas} \\
\midrule
\endhead
\bottomrule
\endfoot

\textbf{¿Qu\'e ordena el primer prove\'ido del juez al recibir la demanda?} &
Ordena el traslado, luego de que el juez revisa si est\'an completos el objeto, el hecho y el derecho invocados. Si falta algo, el prove\'ido ser\'a ``previo'': exige subsanar un requisito antes de continuar el tr\'amite. \\[0.8em]

\textbf{¿Qu\'e es el principio de preclusi\'on?} &
Es el principio por el cual las etapas del proceso se suceden en un orden fijo e irreversible: una vez cerrada una etapa procesal, no se puede volver a ella, y lo que no se hizo en el momento oportuno se pierde como derecho, salvo excepci\'on legal expresa. \\[0.8em]

\textbf{¿Qu\'e regula el art\'iculo 135 del CPCCN?} &
Regula las formas de notificaci\'on de las resoluciones judiciales, en especial la notificaci\'on por c\'edula, que corresponde para el traslado de la demanda y otras resoluciones de particular importancia para el ejercicio del derecho de defensa. \\[0.8em]

\textbf{¿C\'omo se verifica la c\'edula electr\'onica?} &
En los expedientes que tramitan electr\'onicamente (sistema PJN), la c\'edula se genera y consulta dentro del propio expediente digital, disponible desde el inicio de la causa. \\[0.8em]

\textbf{¿Qu\'e es la rebeld\'ia y c\'omo se declara (art\'iculo 59 del CPCCN)?} &
Es la situaci\'on procesal de la parte que, debidamente citada, no contesta la demanda o abandona el proceso. Se declara mediante una resoluci\'on del juez, que debe notificarse para que la parte pueda comparecer en cualquier momento posterior. \\[0.8em]

\textbf{¿Por qu\'e no se aplica la rebeld\'ia en el proceso de alimentos?} &
Porque el proceso de alimentos tiene una funci\'on tuitiva: busca que el alimentante permanezca dentro del proceso y cumpla su obligaci\'on, no que sea excluido de \'el. Por eso puede presentarse en cualquier estado del tr\'amite sin las restricciones propias de la rebeld\'ia. \\[0.8em]

\textbf{¿Qu\'e es la vinculaci\'on provisoria solicitada en el taller pr\'actico?} &
Es una medida urgente pedida dentro de un proceso de familia ante una situaci\'on de temor manifestada por una de las partes, orientada a resolver de manera transitoria el contacto entre las personas involucradas mientras se sustancia el proceso principal. \\[0.8em]

\textbf{¿Qu\'e diferencia hay entre domicilio real y domicilio constituido?} &
El domicilio real es el lugar f\'isico donde efectivamente vive la persona (en el caso trabajado en clase, Moreno 2472, San Miguel, provincia de Buenos Aires); el domicilio constituido es el que se fija a los efectos de las notificaciones dentro del expediente. Ambos deben quedar correctamente identificados, incluyendo la localidad, para que produzcan efectos v\'alidos. \\[0.8em]

\textbf{¿C\'omo se incorpora un pedido nuevo al petitorio de un escrito?} &
No alcanza con mencionarlo en el cuerpo del escrito: debe redactarse con precisi\'on dentro de la secci\'on de petitorio, identificando el punto correspondiente, ya que es esa secci\'on la que delimita exactamente lo que se solicita al juez que resuelva. \\

\midrule
\multicolumn{2}{@{}p{\dimexpr\cornellL+\cornellR+1em\relax}@{}}{%
\textbf{Resumen:} La clase une teor\'ia y pr\'actica en un mismo recorrido: desde que se presenta una demanda hasta que un escrito real llega a la firma. El hilo conductor es el respeto por la forma procesal como garant\'ia del derecho de defensa: el primer prove\'ido del juez, el principio de preclusi\'on, el r\'egimen de notificaciones del art\'iculo 135 del CPCCN y la declaraci\'on de rebeld\'ia del art\'iculo 59 no son formalismos vac\'ios, sino mecanismos que ordenan el proceso y protegen a las partes. La excepci\'on de la rebeld\'ia en el proceso de alimentos muestra que las normas procesales est\'an orientadas por una finalidad de protecci\'on: garantizar que el alimentante permanezca vinculado a la causa. Finalmente, el taller de redacci\'on evidencia que todo ese andamiaje te\'orico se traduce en tareas concretas y urgentes: verificar un domicilio, redactar con precisi\'on un petitorio y coordinar la presentaci\'on de un escrito dentro de los plazos que impone la urgencia del caso. Dominar ambos planos, el conceptual y el operativo, permite ejercer la profesi\'on con responsabilidad frente a las personas reales cuyos derechos est\'an en juego.} \\
\end{longtable}

\end{document}
